%---------------------------------------------------------------------
% UNIT 5 --- MATRICES
% Source: Section 5.8 Self Assessment Questions, pp. 107-108.
%---------------------------------------------------------------------
\chapter{Matrices}

\srcnote{Source: \emph{Business Mathematics} 5405/1429, \S5.8 \saq{Self Assessment Questions}, pp.~107--108.}

\begin{keyidea}[Conformability, and the one rule that is not like arithmetic]
$A_{m\times n}B_{n\times p}$ is defined only when the \emph{inner} dimensions
agree, and the product is $m\times p$. Entry $(i,j)$ of $AB$ is row $i$ of $A$
dotted with column $j$ of $B$.

\medskip
Addition and scalar multiplication behave like arithmetic. Multiplication does
not: in general $AB\neq BA$. The transpose rules follow from this:
\[
(A+B)^{t}=A^{t}+B^{t},\qquad (A^{t})^{t}=A,\qquad
\boxed{(AB)^{t}=B^{t}A^{t}}
\]
--- note the reversal in the last one.
\end{keyidea}

%---------------------------------------------------------------------
\section{Transposes}

\begin{question}{Q.1}
Find the transpose of each of the following matrices:
\[
\text{(i) } A=\mat{-1&0&0\\6&9&-2\\-4&5&1}\qquad
\text{(ii) } B=\mat{6&-6\\4&9\\-2&0}\qquad
\text{(iii) } C=\mat{0&0&2}
\]
\end{question}

\begin{solution}
The transpose turns rows into columns: $(A^{t})_{ij}=A_{ji}$. A matrix of order
$m\times n$ transposes to one of order $n\times m$.

\[
\text{(i) } A^{t}=\mat{-1&6&-4\\0&9&5\\0&-2&1}
\qquad
\text{(ii) } B^{t}=\mat{6&4&-2\\-6&9&0}
\qquad
\text{(iii) } C^{t}=\mat{0\\0\\2}
\]

Orders: $A$ is $3\times3$ so $A^{t}$ is $3\times3$; $B$ is $3\times 2$ so
$B^{t}$ is $2\times 3$; $C$ is a $1\times3$ row vector, so $C^{t}$ is a
$3\times1$ column vector.

\ans{As displayed above. Note that transposing reverses the order:
$3\times2 \to 2\times3$.}
\end{solution}

%---------------------------------------------------------------------
\section{Verifying the double transpose}

\begin{question}{Q.2}
If $A=\mat{-8&0&7\\5&2&-1\\-4&0&3}$, verify that $(A^{t})^{t}=A$.
\end{question}

\begin{solution}
Transpose once --- rows of $A$ become columns:
\[
A^{t}=\mat{-8&5&-4\\0&2&0\\7&-1&3}.
\]
Transpose again --- rows of $A^{t}$ become columns:
\[
(A^{t})^{t}=\mat{-8&0&7\\5&2&-1\\-4&0&3}=A . \qquad\checkmark
\]
Every entry returns to its starting position, because
$\bigl((A^{t})^{t}\bigr)_{ij}=(A^{t})_{ji}=A_{ij}$.

\ans{$(A^{t})^{t}=A$, verified entry by entry.}
\end{solution}

%---------------------------------------------------------------------
\section{Matrix algebra with two $3\times3$ matrices}

\begin{question}{Q.3}
Let $A=\mat{1&2&3\\-1&-2&3\\3&4&5}$ and $B=\mat{4&-1&1\\1&4&3\\3&-3&-2}$.
Find \textbf{(i)} $4A-3B$, \textbf{(ii)} $A+2(B-A)$, \textbf{(iii)}
$2A^{t}+3B^{t}$, \textbf{(iv)} $AB$; and verify that
$(A+B)^{t}=A^{t}+B^{t}$ and $(AB)^{t}=B^{t}A^{t}$.
\end{question}

\begin{solution}
\textbf{(i) $4A-3B$.} Scalar multiplication and subtraction are entrywise:
\[
4A=\mat{4&8&12\\-4&-8&12\\12&16&20},\qquad
3B=\mat{12&-3&3\\3&12&9\\9&-9&-6},
\]
\[
4A-3B=\mat{4-12&8+3&12-3\\-4-3&-8-12&12-9\\12-9&16+9&20+6}
=\mat{-8&11&9\\-7&-20&3\\3&25&26}.
\]

\medskip
\textbf{(ii) $A+2(B-A)$.} Simplify algebraically before touching numbers:
\[
A+2(B-A)=A+2B-2A=2B-A .
\]
\[
2B=\mat{8&-2&2\\2&8&6\\6&-6&-4},\qquad
2B-A=\mat{8-1&-2-2&2-3\\2+1&8+2&6-3\\6-3&-6-4&-4-5}
=\mat{7&-4&-1\\3&10&3\\3&-10&-9}.
\]

\medskip
\textbf{(iii) $2A^{t}+3B^{t}$.}
\[
A^{t}=\mat{1&-1&3\\2&-2&4\\3&3&5},\qquad
B^{t}=\mat{4&1&3\\-1&4&-3\\1&3&-2},
\]
\[
2A^{t}+3B^{t}
=\mat{2&-2&6\\4&-4&8\\6&6&10}+\mat{12&3&9\\-3&12&-9\\3&9&-6}
=\mat{14&1&15\\1&8&-1\\9&15&4}.
\]

\medskip
\textbf{(iv) $AB$.} Row of $A$ against column of $B$. Taking the first row
$(1,\,2,\,3)$ against each column of $B$:
\[
\begin{aligned}
(1)(4)+(2)(1)+(3)(3)&=4+2+9=15,\\
(1)(-1)+(2)(4)+(3)(-3)&=-1+8-9=-2,\\
(1)(1)+(2)(3)+(3)(-2)&=1+6-6=1 .
\end{aligned}
\]
Repeating for rows 2 and 3:
\[
AB=\mat{15&-2&1\\3&-16&-13\\31&-2&5}.
\]

\medskip
\textbf{Verification 1: $(A+B)^{t}=A^{t}+B^{t}$.}
\[
A+B=\mat{5&1&4\\0&2&6\\6&1&3}
\;\Longrightarrow\;
(A+B)^{t}=\mat{5&0&6\\1&2&1\\4&6&3},
\]
\[
A^{t}+B^{t}=\mat{1&-1&3\\2&-2&4\\3&3&5}+\mat{4&1&3\\-1&4&-3\\1&3&-2}
=\mat{5&0&6\\1&2&1\\4&6&3}. \qquad\checkmark
\]

\medskip
\textbf{Verification 2: $(AB)^{t}=B^{t}A^{t}$.}
\[
(AB)^{t}=\mat{15&3&31\\-2&-16&-2\\1&-13&5}.
\]
Computing $B^{t}A^{t}$ --- first row of $B^{t}$ is $(4,1,3)$, first column of
$A^{t}$ is $(1,2,3)^{t}$, giving $4+2+9=15$, and so on:
\[
B^{t}A^{t}=\mat{15&3&31\\-2&-16&-2\\1&-13&5}. \qquad\checkmark
\]
The order genuinely matters: $A^{t}B^{t}$ is a different matrix entirely.

\ans{(i) $\mat{-8&11&9\\-7&-20&3\\3&25&26}$ \quad
(ii) $\mat{7&-4&-1\\3&10&3\\3&-10&-9}$ \quad
(iii) $\mat{14&1&15\\1&8&-1\\9&15&4}$ \quad
(iv) $AB=\mat{15&-2&1\\3&-16&-13\\31&-2&5}$;
both transpose identities verified.}
\end{solution}

\begin{pitfall}
$(AB)^{t}=B^{t}A^{t}$, \textbf{not} $A^{t}B^{t}$. The reversal is not a
convention --- it is forced by conformability, and writing it the wrong way
round gives a matrix that is usually not even the right shape.
\end{pitfall}

%---------------------------------------------------------------------
\section{A commuting-matrices proof}

\begin{question}{Q.4}
Prove that if $AB=BA$, then $(A+B)(A-B)=A^{2}-B^{2}$.
\end{question}

\begin{solution}
Expand the left side using the distributive law, which \emph{does} hold for
matrices. What does not hold in general is commutativity, so keep every product
in the order it arises:
\[
(A+B)(A-B)=A(A-B)+B(A-B)=A^{2}-AB+BA-B^{2}.
\]
So far this is true for any $A$ and $B$. Now apply the hypothesis $AB=BA$,
which lets the middle two terms cancel:
\[
-AB+BA=-AB+AB=0 .
\]
Therefore
\[
(A+B)(A-B)=A^{2}-B^{2}. \qquad\blacksquare
\]

\medskip
\textbf{The converse is worth noting.} Running the argument backwards, if
$(A+B)(A-B)=A^{2}-B^{2}$ then $BA-AB=0$, so $A$ and $B$ must commute. The
familiar difference-of-squares identity holds for matrices \emph{exactly} when
$AB=BA$ --- which is why the hypothesis is stated.

\ans{Expanding gives $A^{2}-AB+BA-B^{2}$; the hypothesis $AB=BA$ makes
$-AB+BA=0$, leaving $A^{2}-B^{2}$.}
\end{solution}

%---------------------------------------------------------------------
\section{Valuing an inventory by matrix multiplication}

\begin{question}{Q.5}
A pet store has 6 kittens, 8 parrots and 12 puppies. If each kitten costs \$28,
each parrot \$30 and each puppy \$40, use matrix multiplication to find the
total value of the store's inventory.
\end{question}

\begin{solution}
Write the quantities as a $1\times3$ row and the prices as a $3\times1$ column.
The inner dimensions agree ($3=3$), so the product is $1\times1$ --- a single
number, which is exactly what a total value should be.
\[
Q=\mat{6&8&12},\qquad P=\mat{28\\30\\40},
\]
\[
QP=\mat{6&8&12}\mat{28\\30\\40}
=(6)(28)+(8)(30)+(12)(40)
\]
\[
=168+240+480=\mat{888}.
\]

\ans{The inventory is worth \textbf{\$888}.}
\end{solution}

\begin{pitfall}
Order the factors so the inner dimensions match. $PQ$ with $P$ as $3\times1$
and $Q$ as $1\times3$ \emph{is} defined, but produces a $3\times3$ matrix of
every price-times-quantity cross-product --- not a total.
\end{pitfall}

%---------------------------------------------------------------------
\section{Match-day revenue}

\begin{question}{Q.6}
Attendance for the first three football games of the season is given below.
Adult tickets sold for \$3.00 and student tickets for \$1.50. Use matrix
multiplication to find the revenue for each game.

\medskip
\centering\small\sffamily
\begin{tabular}{@{}l C{2.2cm} C{2.2cm}@{}}
\toprule
 & \textbf{Adults} & \textbf{Students} \\
\midrule
Game 1 & 340 & 180 \\
Game 2 & 250 & 195 \\
Game 3 & 300 & 220 \\
\bottomrule
\end{tabular}
\end{question}

\begin{solution}
Attendance is a $3\times2$ matrix (one row per game), prices a $2\times1$
column. Inner dimensions agree, so the product is $3\times1$ --- one revenue
figure per game.
\[
A=\mat{340&180\\250&195\\300&220},\qquad
P=\mat{3.00\\1.50},
\]
\[
AP=\mat{340(3.00)+180(1.50)\\250(3.00)+195(1.50)\\300(3.00)+220(1.50)}
=\mat{1020+270\\750+292.50\\900+330}
=\mat{1290.00\\1042.50\\1230.00}.
\]

\begin{center}
\small\sffamily
\begin{tabular}{@{}l C{2.2cm} C{2.2cm} C{2.6cm}@{}}
\toprule
 & \textbf{Adults @ \$3} & \textbf{Students @ \$1.50} & \textbf{Revenue} \\
\midrule
Game 1 & \$1{,}020 & \$270.00   & \textbf{\$1{,}290.00} \\
Game 2 & \$750     & \$292.50   & \textbf{\$1{,}042.50} \\
Game 3 & \$900     & \$330.00   & \textbf{\$1{,}230.00} \\
\midrule
\textbf{Season} & & & \textbf{\$3{,}562.50} \\
\bottomrule
\end{tabular}
\end{center}

\ans{Game 1: \$1{,}290.00; Game 2: \$1{,}042.50; Game 3: \$1{,}230.00.
Total for the three games: \$3{,}562.50.}
\end{solution}

%---------------------------------------------------------------------
\section{A scalar for a price rise}

\begin{question}{Q.7}
Let $A=\mat{a_1&a_2&a_3}$ represent the prices of products $X$, $Y$, $Z$. If
prices are to be increased by 15\%, by what scalar can $A$ be multiplied to
obtain the new prices?
\end{question}

\begin{solution}
A 15\% increase means each new price is the old price \emph{plus} 15\% of it:
\[
a_i^{\text{new}}=a_i+0.15a_i=1.15\,a_i .
\]
Because the same factor applies to every entry, this is exactly scalar
multiplication:
\[
A^{\text{new}}=1.15\,A=\mat{1.15a_1 & 1.15a_2 & 1.15a_3}.
\]

\ans{Multiply by the scalar $\mathbf{1.15}$ (equivalently $\tfrac{23}{20}$).}
\end{solution}

\begin{pitfall}
The scalar is $1.15$, not $0.15$. Multiplying by $0.15$ gives the size of the
\emph{increase} alone and discards the original price.
\end{pitfall}

%---------------------------------------------------------------------
\section{Burger Barn: three branches, three products}

\begin{question}{Q.8}
Burger Barn's three locations sell fries, burgers and soft drinks. Branch I
sells 500 orders of fries, 150 burgers and 300 soft drinks daily. Branches II
and III sell 900 and 700 orders of fries respectively per day. Branch II sells
440 burgers and 830 soft drinks, while Branch III sells 580 burgers and 1200
soft drinks each day.
\textbf{(a)} Write a $3\times3$ matrix representing daily sales for all
locations.
\textbf{(b)} If fries cost \$0.92 per order, burgers \$1.54 each and soft
drinks \$0.58 each, write a $1\times3$ matrix representing the prices.
\textbf{(c)} What is the total daily income from all the locations?
\end{question}

\begin{solution}
\textbf{(a) The sales matrix.} Rows are branches, columns are products:
\[
S=\begin{array}{c@{\;}c}
 & \begin{array}{ccc}\text{\scriptsize fries}&\text{\scriptsize burgers}&\text{\scriptsize drinks}\end{array}\\
\begin{array}{c}\text{\scriptsize I}\\ \text{\scriptsize II}\\ \text{\scriptsize III}\end{array} &
\mat{500&150&300\\900&440&830\\700&580&1200}
\end{array}
\]

\medskip
\textbf{(b) The price matrix.} As given, a $1\times3$ row:
\[
P=\mat{0.92&1.54&0.58}.
\]

\medskip
\textbf{(c) Total daily income.} $S$ is $3\times3$ and $P$ is $1\times3$, so
$SP$ is \emph{not} conformable. Transpose the prices into a column first, and
the inner dimensions agree:
\[
SP^{t}=\mat{500&150&300\\900&440&830\\700&580&1200}\mat{0.92\\1.54\\0.58}.
\]
Branch by branch:
\[
\begin{aligned}
\text{I}  &: 500(0.92)+150(1.54)+300(0.58)=460+231+174=865.00,\\
\text{II} &: 900(0.92)+440(1.54)+830(0.58)=828+677.60+481.40=1987.00,\\
\text{III}&: 700(0.92)+580(1.54)+1200(0.58)=644+893.20+696=2233.20 .
\end{aligned}
\]
\[
SP^{t}=\mat{865.00\\1987.00\\2233.20}.
\]
Total across all branches:
\[
865.00+1987.00+2233.20=5085.20 .
\]

\begin{center}
\small\sffamily
\begin{tabular}{@{}l C{2.2cm} C{2.2cm} C{2.2cm} C{2.6cm}@{}}
\toprule
\textbf{Branch} & \textbf{Fries} & \textbf{Burgers} & \textbf{Drinks} & \textbf{Daily income} \\
\midrule
I   & 500 & 150 & 300  & \$865.00 \\
II  & 900 & 440 & 830  & \$1{,}987.00 \\
III & 700 & 580 & 1200 & \$2{,}233.20 \\
\midrule
\multicolumn{4}{@{}l}{\textbf{Total}} & \textbf{\$5{,}085.20} \\
\bottomrule
\end{tabular}
\end{center}

\ans{(a) $S=\mat{500&150&300\\900&440&830\\700&580&1200}$ \quad
(b) $P=\mat{0.92&1.54&0.58}$ \quad
(c) Branch incomes \$865.00, \$1{,}987.00 and \$2{,}233.20, giving a total daily
income of \textbf{\$5{,}085.20}.}
\end{solution}

\begin{pitfall}
The question deliberately asks for $P$ as a \emph{row}, then asks for a total
that requires it as a \emph{column}. Say so in your working: ``$SP$ is not
conformable, so compute $SP^{t}$.'' Examiners award that observation.
\end{pitfall}
